\newif\ifglobal

%%%%%%%%%%%%%%%%%%%%%%%
%%% Compile to view %%%
%%%%%%%%%%%%%%%%%%%%%%%

%%%%%%%%%%%%%%%%%%%%%%%%%%%%%%%%%%%%%%%%%%%%%%%%%%%
%%% Readme:                                     %%%
%%% https://www.overleaf.com/read/bwdjvfjksvjy/ %%%
%%%%%%%%%%%%%%%%%%%%%%%%%%%%%%%%%%%%%%%%%%%%%%%%%%%

% >>> M?: marks a module setting number or important notice

% >>> M4: Set {./relative-path} to external files for this module <<<
\def \modulefiles {./19-DBGASSIST}
% To add an external file, use:
% (Replace '---' with actual filename)
% '\input{\modulefiles/tables/---.tex}' for tables
% '\includegraphics{\modulefiles/figures/---}' for figures
% '\subfile{\modulefiles/---}' for register files

% >>> M1: This module is in global project? <<<
% 'YES' - uncomment; 'NO' - comment out
\globaltrue

\ifglobal
    \documentclass[main\_\_EN.tex]{subfiles}

\else
    \documentclass[openany, 10pt]{book}

    % >>> M2: In ./00-shared/preamble-trm-module, also find
    %         settings for visibility of todo notes, watermark <<<
    \usepackage{preamble-trm-module}

    % >>> M3: Temporary labels in English,
    %         you can add your own labels there <<<
    \myexternaldocument{./00-shared/temp-labels-en}

\fi %\ifglobal


\setcounter{secnumdepth}{4}


% Define formatting of parts and chapters in text
\titleformat{\chapter}
  [display]
  {\LARGE\bfseries}
  {\color{AllportsColor}Chapter \thechapter}
  {1em}
  {\color{AllportsColor}}
\titlespacing*{\chapter}{0pt}{0pt}{20pt}

% Define headers
\usepackage{fancyhdr}
\usepackage{lipsum}
\renewcommand{\chaptermark}[1]{\markboth{Chapter\, \thechapter\, \ #1}{}}
\fancypagestyle{plain}{%
  \fancyhead{}%
  %\fancyhead[R]{\Navigation}
  \fancyhead[L]{%
    \nouppercase{%
      \ %
      \normalsize\leftmark
    }%
  }%
}
\pagestyle{plain}


% >>> M5: If this module needs any updates in
% './00-shared/preamble-trm-module.sty'
% (include additional package, etc.), contact your module PM
% or create the file './00-shared/changelog.tex' make the
% necessary updates yourself and document them in this file <<<

% Make text left-aligned and allow latex to break pages naturally, comment out for full justification
\raggedright
\raggedbottom


\begin{document}


% Sets language into which pre-defined bits of text are translated
\selectlanguage{english}

% Do not delete this block
\ifglobal
% Add the following front matter if
% this module is outside of global project
\else
    \InputIfFileExists{readme}{}{}
    \listoftodos
    {\let\clearpage\relax \listofdonetodos}
    \tableofcontents
    %\listoftables
    %\listoffigures
\fi


%%%%%%%%%%%%%%%%%%%%%%%%%
%%% TRM Module Begins %%%
%%%%%%%%%%%%%%%%%%%%%%%%%

\chapter{Debug Assistant}
\label{mod:dbgassist}
\hypertarget{dbgassist}{}

\section{Overview}

Debug Assistant is an auxiliary module that features a set of functions to help locate bugs and issues during software debugging.

\section{Features}
\label{sec:dbgassist-features}

The Debug Assistant module has the following features:

\begin{itemize}
    \item \textbf{Region read/write monitoring}: Monitors whether the High-Performance CPU (HP CPU) bus reads from or writes to a specified memory address space. A detected read or write in the monitored address space triggers an interrupt.
    \item \textbf{Stack pointer (SP) monitoring}: Monitors whether the SP exceeds the specified address space. A bounds violation triggers an interrupt.
    \item \textbf{Program counter (PC) logging}: Records the PC value. The last PC value at the most recent reset of the HP CPU can be retrieved.
    \item \textbf{Bus access logging}: Records information about bus access. When the HP CPU or the Direct Memory Access controller (DMA) writes a specified value, the Debug Assistant module records the data type, address of the write operation, and the PC value when the write is performed by the HP CPU, then stores this information in the HP SRAM.
\end{itemize}

\section{Functional Description}
\label{sec:dbgassist-func-descr}

\subsection{Region Read/Write Monitoring}

The Debug Assistant module monitors reads and writes performed by the HP CPU bus in a specified address space, i.e., memory region. Whenever the bus reads from or writes to this address space, an interrupt is triggered. Refer to Table \ref{tab:riscvcpu-mem-blocks} in Chapter \ref{mod:riscvcpu} \textit{\nameref{mod:riscvcpu}}.

In this chapter, when the HP CPU bus accesses the data space, it is referred to as the \textbf{Data bus}. When the HP CPU bus accesses the AHB peripheral space, it is referred to as the \textbf{Peripheral bus}. The Debug Assistant module can simultaneously monitor two memory regions (HP CPU region 0 and HP CPU region 1) on the Data bus and the Peripheral bus. The regions should be defined based on the developer's needs. 

The supported region read/write monitoring modes are:

\begin{itemize}
    \item Monitoring read/write operations on the Data bus:
        \begin{itemize}
            \item HP CPU Data bus reads from HP CPU region 0.
            \item HP CPU Data bus writes to HP CPU region 0.
            \item HP CPU Data bus reads from HP CPU region 1.
            \item HP CPU Data bus writes to HP CPU region 1.
        \end{itemize}    
    \item Monitoring read/write operations on the Peripheral bus:
        \begin{itemize}
            \item HP CPU Peripheral bus reads from HP CPU region 0.
            \item HP CPU Peripheral bus writes to HP CPU region 0.
            \item HP CPU Peripheral bus reads from HP CPU region 1.
            \item HP CPU Peripheral bus writes to HP CPU region 1.
        \end{itemize}
\end{itemize}

\subsection{SP Monitoring}

The Debug Assistant module monitors the SP to prevent stack overflow or erroneous push/pop operations in the HP CPU. When the HP CPU SP exceeds the monitored region's lower or upper bounds, the module records the PC and generates an interrupt. The bounds are configured by software. 

The supported SP monitoring modes are:

\begin{itemize}
    \item The HP CPU SP exceeds the upper bound of the monitored region.
    \item The HP CPU SP falls below the lower bound of the monitored region.
\end{itemize}

\subsection{PC Logging}

Software developers may need to identify the PC value at the last reset of the HP CPU. For instance, when the program is stuck and requires a reset, the developer may want to know where the program got stuck to debug. The Debug Assistant module records the PC at the last reset of the HP CPU, which can then be read for debugging.

\subsection{CPU/DMA Bus Access Logging}

The Debug Assistant module records write operations of the HP CPU Data bus and the DMA bus in real time. When a write operation occurs or a specific value is written to a specified address space, the module records the bus type, address, PC (only when the write is performed by the HP CPU), and other information. It then stores the data in the HP SRAM in a specific format. The specified address range must fall within the permitted memory regions of HP SRAM or external RAM.

%%%%%%%%%%%%%%%%%%
%%% Interrupts %%%
%%%%%%%%%%%%%%%%%%

\section{Interrupts}
\label{sec:dbgassist-interrupts}

The Debug Assistant can generate the BUS\_MONITOR\_INT interrupt signal that will be sent to the \hyperref[mod:intmtrx]{Interrupt Matrix}. The following interrupt sources can generate this interrupt signal.

\begin{itemize}
    \item \label{int:BUSMONITORCORE0AREADRAM00RDINT}BUS\_MONITOR\_CORE\_0\_AREA\_DRAM0\_0\_RD\_INT: Triggered when the HP CPU Data bus reads from HP CPU region 0.
    \item \label{int:BUSMONITORCORE0AREADRAM00WRINT}BUS\_MONITOR\_CORE\_0\_AREA\_DRAM0\_0\_WR\_INT: Triggered when the HP CPU Data bus writes to HP CPU region 0.
    \item \label{int:BUSMONITORCORE0AREADRAM01RDINT}BUS\_MONITOR\_CORE\_0\_AREA\_DRAM0\_1\_RD\_INT: Triggered when the HP CPU Data bus reads from HP CPU region 1.
    \item \label{int:BUSMONITORCORE0AREADRAM01WRINT}BUS\_MONITOR\_CORE\_0\_AREA\_DRAM0\_1\_WR\_INT: Triggered when the HP CPU Data bus writes to HP CPU region 1.
    \item \label{int:BUSMONITORCORE0AREAPIF0RDINT}BUS\_MONITOR\_CORE\_0\_AREA\_PIF\_0\_RD\_INT: Triggered when the HP CPU Peripheral bus reads from HP CPU region 0.
    \item \label{int:BUSMONITORCORE0AREAPIF0WRINT}BUS\_MONITOR\_CORE\_0\_AREA\_PIF\_0\_WR\_INT: Triggered when the HP CPU Peripheral bus writes to HP CPU region 0.
    \item \label{int:BUSMONITORCORE0AREAPIF1RDINT}BUS\_MONITOR\_CORE\_0\_AREA\_PIF\_1\_RD\_INT: Triggered when the HP CPU Peripheral bus reads from HP CPU region 1.
    \item \label{int:BUSMONITORCORE0AREAPIF1WRINT}BUS\_MONITOR\_CORE\_0\_AREA\_PIF\_1\_WR\_INT: Triggered when the HP CPU Peripheral bus writes to HP CPU region 1.
    \item \label{int:BUSMONITORCORE0SPSPILLMININT}BUS\_MONITOR\_CORE\_0\_SP\_SPILL\_MIN\_INT: Triggered when the HP CPU SP falls below the lower bound of the monitored region.
    \item \label{int:BUSMONITORCORE0SPSPILLMAXINT}BUS\_MONITOR\_CORE\_0\_SP\_SPILL\_MAX\_INT: Triggered when the HP CPU SP exceeds the upper bound of the monitored region.
\end{itemize}

\begin{tiplisting}
For definitions of \textit{interrupt}, \textit{interrupt signal}, \textit{interrupt source}, and their correlations, please refer to Chapter \ref{mod:intmtrx} \textit{\nameref{mod:intmtrx}} > Section \ref{sec:intmtx-int-terminology} \textit{\nameref{sec:intmtx-int-terminology}}.
\end{tiplisting}

Each interrupt source can be configured by a common set of registers that are described in Section \hyperref[interrupt-config-registers]{\textit{Interrupt Configuration Registers}}. The specific registers can be found in Section \ref{sec:dbgassist-reg-summ} \textit{\nameref{sec:dbgassist-reg-summ}}.

\section{Programming Procedures}
\label{sec:dbgassist-prog-procedures}

\subsection{Region Monitoring and SP Monitoring Configuration}

The configuration process for region monitoring and SP monitoring is as follows:

\begin{enumerate}
    \item Configure the monitored region and SP bounds.
        \begin{itemize}
            \item Set the HP CPU Data bus region 0 with \hyperref[regdesc:BUSMONITORCORE0AREADRAM00MINREG]{BUS\_MONITOR\_CORE\_0\_AREA\_DRAM0\_0\_MIN\_REG} and \hyperref[regdesc:BUSMONITORCORE0AREADRAM00MAXREG]{BUS\_MONITOR\_CORE\_0\_AREA\_DRAM0\_0\_MAX\_REG}. 
            \item Set the HP CPU Data bus region 1 with \hyperref[regdesc:BUSMONITORCORE0AREADRAM01MINREG]{BUS\_MONITOR\_CORE\_0\_AREA\_DRAM0\_1\_MIN\_REG} and \hyperref[regdesc:BUSMONITORCORE0AREADRAM01MAXREG]{BUS\_MONITOR\_CORE\_0\_AREA\_DRAM0\_1\_MAX\_REG}.
            \item Set the HP CPU Peripheral bus region 0 with \hyperref[regdesc:BUSMONITORCORE0AREAPIF0MINREG]{BUS\_MONITOR\_CORE\_0\_AREA\_PIF\_0\_MIN\_REG} and \hyperref[regdesc:BUSMONITORCORE0AREAPIF0MAXREG]{BUS\_MONITOR\_CORE\_0\_AREA\_PIF\_0\_MAX\_REG}.
            \item Set the HP CPU Peripheral bus region 1 with \hyperref[regdesc:BUSMONITORCORE0AREAPIF1MINREG]{BUS\_MONITOR\_CORE\_0\_AREA\_PIF\_1\_MIN\_REG} and \hyperref[regdesc:BUSMONITORCORE0AREAPIF1MAXREG]{BUS\_MONITOR\_CORE\_0\_AREA\_PIF\_1\_MAX\_REG}.
            \item Set the HP CPU SP bounds with \hyperref[regdesc:BUSMONITORCORE0SPMINREG]{BUS\_MONITOR\_CORE\_0\_SP\_MIN\_REG} and \hyperref[regdesc:BUSMONITORCORE0SPMAXREG]{BUS\_MONITOR\_CORE\_0\_SP\_MAX\_REG}.
        \end{itemize}
    \item Set \hyperref[regdesc:BUSMONITORCORE0INTRENAREG]{BUS\_MONITOR\_CORE\_0\_INTR\_ENA\_REG} to enable interrupts for monitoring mode.
    \item Set \hyperref[regdesc:BUSMONITORCORE0MONTRENAREG]{BUS\_MONITOR\_CORE\_0\_MONTR\_ENA\_REG} to enable the monitoring mode(s). You can enable multiple monitoring modes simultaneously.
    \item Read \hyperref[regdesc:BUSMONITORCORE0INTRRAWREG]{BUS\_MONITOR\_CORE\_0\_INTR\_RAW\_REG} to obtain the raw interrupt status of a monitoring mode. 
    \item Set \hyperref[regdesc:BUSMONITORCORE0INTRCLRREG]{BUS\_MONITOR\_CORE\_0\_INTR\_CLR\_REG} to clear the interrupt for a monitoring mode.
\end{enumerate}

For example, if the Debug Assistant module needs to monitor whether the HP CPU Data bus has written to the [A \verb+~+ B] address space, the user can enable monitoring in either the HP CPU Data bus region 0 or region 1. The following configuration process is based on region 0:

\begin{enumerate}
    \item Set \hyperref[fielddesc:BUSMONITORCORE0RCDPDEBUGEN]{BUS\_MONITOR\_CORE\_0\_RCD\_PDEBUGEN} to 1 to enable the HP CPU to update the PC signals to the Debug Assistant module.
    \item Set \hyperref[regdesc:BUSMONITORCORE0AREADRAM00MINREG]{BUS\_MONITOR\_CORE\_0\_AREA\_DRAM0\_0\_MIN\_REG} to Address A.
    \item Set \hyperref[regdesc:BUSMONITORCORE0AREADRAM00MAXREG]{BUS\_MONITOR\_CORE\_0\_AREA\_DRAM0\_0\_MAX\_REG} to Address B.
    \item Set \hyperref[regdesc:BUSMONITORCORE0INTRENAREG]{BUS\_MONITOR\_CORE\_0\_INTR\_ENA\_REG} bit[1] to enable the interrupt for write operations by the HP CPU Data bus in region 0.
    \item Set \hyperref[regdesc:BUSMONITORCORE0MONTRENAREG]{BUS\_MONITOR\_CORE\_0\_MONTR\_ENA\_REG} bit[1] to enable monitoring of write operations by the HP CPU Data bus in region 0.
    \item Configure the interrupt matrix to map BUS\_MONITOR\_INTR into the HP CPU interrupt (refer to Chapter \ref{mod:intmtrx} \textit{\nameref{mod:intmtrx}}).
    \item After the interrupt is triggered:
        \begin{itemize}
            \item Read \hyperref[regdesc:BUSMONITORCORE0INTRRAWREG]{BUS\_MONITOR\_CORE\_0\_INTR\_RAW\_REG} to identify the interrupt source.
            \item If the interrupt is triggered by the HP CPU region monitoring, read \hyperref[fielddesc:BUSMONITORCORE0AREAPC]{BUS\_MONITOR\_CORE\_0\_AREA\_PC} for the HP CPU PC value and \hyperref[fielddesc:BUSMONITORCORE0AREASP]{BUS\_MONITOR\_CORE\_0\_AREA\_SP} for the HP CPU SP.
            \item If the interrupt is triggered by SP monitoring, read \hyperref[fielddesc:BUSMONITORCORE0SPPC]{BUS\_MONITOR\_CORE\_0\_SP\_PC} for the HP CPU PC value.
            \item Write 1 to the corresponding bit of \hyperref[regdesc:BUSMONITORCORE0INTRCLRREG]{BUS\_MONITOR\_CORE\_0\_INTR\_CLR\_REG} to clear the interrupt.
        \end{itemize}
\end{enumerate}

\subsection{PC Logging Configuration}

Configure \hyperref[fielddesc:BUSMONITORCORE0RCDPDEBUGEN]{BUS\_MONITOR\_CORE\_0\_RCD\_PDEBUGEN} to 1 to enable the HP CPU to update the PC signals to the Debug Assistant module. If \hyperref[fielddesc:BUSMONITORCORE0RCDRECORDEN]{BUS\_MONITOR\_CORE\_0\_RCD\_RECORDEN} is also set to 1, \hyperref[regdesc:BUSMONITORCORE0RCDPDEBUGPCREG]{BUS\_MONITOR\_CORE\_0\_RCD\_PDEBUGPC\_REG} will record the HP CPU's PC value, and \hyperref[regdesc:BUSMONITORCORE0RCDPDEBUGSPREG]{BUS\_MONITOR\_CORE\_0\_RCD\_PDEBUGSP\_REG} will record the HP CPU's SP value. Otherwise, these two registers will retain their original values.

When the HP CPU resets, \hyperref[regdesc:BUSMONITORCORE0RCDENREG]{BUS\_MONITOR\_CORE\_0\_RCD\_EN\_REG} will reset, while \hyperref[regdesc:BUSMONITORCORE0RCDPDEBUGPCREG]{BUS\_MONITOR\_CORE\_0\_RCD\_PDEBUGPC\_REG} and \hyperref[regdesc:BUSMONITORCORE0RCDPDEBUGSPREG]{BUS\_MONITOR\_CORE\_0\_RCD\_PDEBUGSP\_REG} will not. Therefore, these two registers will retain the PC and SP values during the HP CPU reset.

\subsection{CPU/DMA Bus Access Logging Configuration}

\subsubsection{Bus Access Logging 1 Configuration}
\label{sec:hp-cpu-access-mon}

Bus access logging 1 includes HP CPU bus access logging (HP SRAM and external RAM monitoring) and DMA bus access logging (HP SRAM monitoring). The configuration process is described below.

\begin{enumerate}
    \item Configure the monitored address space: set \hyperref[regdesc:MEMMONITORLOGMINREG]{MEM\_MONITOR\_LOG\_MIN\_REG} and \hyperref[regdesc:MEMMONITORLOGMAXREG]{MEM\_MONITOR\_LOG\_MAX\_REG} to specify the monitored address space. For HP CPU, the monitored address space should be in the range of HP SRAM (0x4080\_0000 \verb+~+ 0x4084\_FFFF) or external RAM (0x4200\_0000 \verb+~+ 0x43FF\_FFFF). For DMA, the monitored address space should be in the range of HP SRAM (0x4080\_0000 \verb+~+ 0x4084\_FFFF).
    \item Configure the monitoring mode with \hyperref[fielddesc:MEMMONITORLOGMODE]{MEM\_MONITOR\_LOG\_MODE}:
        \begin{itemize}
            \item Write monitoring (detects bus write operations)
            \item Word monitoring (detects writes of a specific word)
            \item Halfword monitoring (detects writes of a specific halfword)
            \item Byte monitoring (detects writes of a specific byte)
        \end{itemize}
    \item Configure the specific values to monitor:
        \begin{itemize}
            \item In word monitoring mode, \hyperref[regdesc:MEMMONITORLOGCHECKDATAREG]{MEM\_MONITOR\_LOG\_CHECK\_DATA\_REG} specifies the monitored word.
            \item In halfword monitoring mode, \hyperref[regdesc:MEMMONITORLOGCHECKDATAREG]{MEM\_MONITOR\_LOG\_CHECK\_DATA\_REG}[15:0] specifies the monitored halfword.
            \item In byte monitoring mode, \hyperref[regdesc:MEMMONITORLOGCHECKDATAREG]{MEM\_MONITOR\_LOG\_CHECK\_DATA\_REG}[7:0] specifies the monitored byte.
            \item \hyperref[regdesc:MEMMONITORLOGDATAMASKREG]{MEM\_MONITOR\_LOG\_DATA\_MASK\_REG} masks the byte(s) specified in \hyperref[regdesc:MEMMONITORLOGCHECKDATAREG]{MEM\_MONITOR\_LOG\_CHECK\_DATA\_REG}. A masked byte can be any value. For example, in word monitoring, if \hyperref[regdesc:MEMMONITORLOGCHECKDATAREG]{MEM\_MONITOR\_LOG\_CHECK\_DATA\_REG} is set to 0x01020304 and \hyperref[regdesc:MEMMONITORLOGDATAMASKREG]{MEM\_MONITOR\_LOG\_DATA\_MASK\_REG} is set to 0x1, then any writes matching the 0x010203XX pattern will be recorded.
        \end{itemize}
    \item Configure the storage space for recorded data:
        \begin{itemize}
            \item \hyperref[regdesc:MEMMONITORLOGMEMSTARTREG]{MEM\_MONITOR\_LOG\_MEM\_START\_REG} and \hyperref[regdesc:MEMMONITORLOGMEMENDREG]{MEM\_MONITOR\_LOG\_MEM\_END\_REG} specify the storage space for recorded data. The storage space must be in the range of 0x4080\_0000 \verb+~+ 0x4084\_FFFF.
            \item Set \hyperref[regdesc:MEMMONITORLOGMEMADDRUPDATEREG]{MEM\_MONITOR\_LOG\_MEM\_ADDR\_UPDATE\_REG} to update the value in \hyperref[regdesc:MEMMONITORLOGMEMCURRENTADDRREG]{MEM\_MONITOR\_LOG\_MEM\_CURRENT\_ADDR\_REG} to \hyperref[regdesc:MEMMONITORLOGMEMSTARTREG]{MEM\_MONITOR\_LOG\_MEM\_START\_REG}.
            \item Configure the Debug Assistant module's permission to access HP SRAM. The Debug Assistant module can access HP SRAM only when the access permission is enabled. For more information, refer to Chapter \ref{mod:permctrl} \textit{\nameref{mod:permctrl}}.
        \end{itemize}
    \item \label{other:dbgassist-item-log-config-step5} Configure the writing mode for recorded data for loop mode or non-loop mode:
        \begin{itemize}
            \item In loop mode, writing to the specified address space occurs in loops. When writing reaches the end address, it returns to the starting address and continues, overwriting previously recorded data. Set \hyperref[fielddesc:MEMMONITORLOGMEMLOOPENABLE]{MEM\_MONITOR\_LOG\_MEM\_LOOP\_ENABLE} to 1 to enable loop mode.\\
            For example, if there are 10 write operations (1 \verb+~+ 10) to address space 0 \verb+~+ 4 during bus access, after the 5th operation writes to address 4, the 6th operation will start writing from address 0. The 6th to 10th operations will overwrite the data written by the 1st to 5th operations.
            \item In non-loop mode, when writing reaches the end address, it stops at the end address and discards the remaining data. The previously recorded data will not be overwritten. Clear \hyperref[fielddesc:MEMMONITORLOGMEMLOOPENABLE]{MEM\_MONITOR\_LOG\_MEM\_LOOP\_ENABLE} to enable non-loop mode.\\
            For example, if there are 10 write operations (1 \verb+~+ 10) to address space 0 \verb+~+ 4 during bus access, after the 5th operation writes to address 4, the 6th to 10th write operations will stop at address 4 and will not be performed. Therefore, address 0 \verb+~+ 4 will store the values written by the 1st to 5th operations, while the values of the 6th to 10th operations will be discarded.
        \end{itemize}
    \item Configure bus enable registers.
        \begin{itemize}
            \item Enable HP CPU and DMA bus access logging respectively with \hyperref[fielddesc:MEMMONITORLOGCOREENA]{MEM\_MONITOR\_LOG\_CORE\_ENA} and \hyperref[fielddesc:MEMMONITORLOGDMA0ENA]{MEM\_MONITOR\_LOG\_DMA\_0\_ENA}. They can be enabled at the same time.
        \end{itemize}
\end{enumerate}

The Debug Assistant module first writes the recorded data to internal buffer A, then fetches the data from the buffer and writes it to the configured memory space. When monitored behaviors are triggered continuously, the generated recording packets may fill the buffer, preventing it from accepting new packets. In such cases, the module discards incoming packets and buffers a LOST packet instead before the buffer reaches full capacity. The DMA LOST packet contains information about the bus type of the discarded packets but not the number of discarded packets.

When bus access logging is finished, the recorded data can be read from memory for decoding. The recorded data consists of three packet formats: HP CPU packet (for the HP CPU Data bus), DMA packet (for the DMA Data bus), and LOST packet. The packet formats are shown in Tables \ref{tab:dbgassist-hpcpu-package}, \ref{tab:dbgassist-dma-package}, and \ref{tab:dbgassist-lost-package}.

\begin{longtable}[c]{|L{3cm}|L{3cm}|L{3cm}|L{3cm}|L{3cm}|}
    \caption{HP CPU Packet Format}\label{tab:dbgassist-hpcpu-package}
    \\\hline
    \rowcolor{lightgray}
    \textbf{Bit[63:33]} & \textbf{Bit[32]} & \textbf{Bit[31:3]} & \textbf{Bit[2:1]} & \textbf{Bit[0]} \\\hline
    pc\_offset & anchored(1) & addr\_offset & format & anchored(0) \\\hline
\end{longtable}

\begin{longtable}[c]{|L{2cm}|L{2cm}|L{2cm}|L{2cm}|}
    \caption{DMA Packet Format}\label{tab:dbgassist-dma-package}
    \\\hline
    \rowcolor{lightgray}
    \textbf{Bit[31:8]} & \textbf{Bit[7:3]} & \textbf{Bit[2:1]} & \textbf{Bit[0]} \\\hline
    addr\_offset & dma\_source & format & anchored(0) \\\hline
\end{longtable}

\begin{longtable}[c]{|L{2cm}|L{2cm}|L{2cm}|L{2cm}|}
    \caption{LOST Packet Format}\label{tab:dbgassist-lost-package}
    \\\hline
    \rowcolor{lightgray}
    \textbf{Bit[31:7]} & \textbf{Bit[6:3]} & \textbf{Bit[2:1]} & \textbf{Bit[0]} \\\hline
    reserved & lost\_source & format & anchored(0) \\\hline
\end{longtable}

The data packet formats show that the HP CPU packet is 64 bits, while the other packets are 32 bits each. These packets have the following fields:

\begin{itemize}
    \item \textbf{format} – indicates the source of the packet type:
    \begin{itemize}
        \item 0: HP CPU packet
        \item 1: DMA packet
        \item 2: Reserved
        \item 3: LOST packet
    \end{itemize}
    \item \textbf{pc\_offset} - indicates the offset of the HP CPU PC register at the time of access. Actual PC = pc\_offset + 0x0000\_0000.
    \item \textbf{addr\_offset} - the address offset of a write operation. Actual address = addr\_offset + \{\hyperref[regdesc:MEMMONITORLOGMINREG]{MEM\_MONITOR\_LOG\_MIN\_REG}[31:2], 2'b0\}.
    \item \textbf{dma\_source} - the source of DMA access. Details can be found in Table \ref{tab:assistdbg-dma-source}, where sources corresponding to values 16 \verb+~+ 31 are detailed in \ref{mod:dma} \textit{\nameref{mod:dma}}.
    \item \textbf{lost\_source} - indicates which packets were discarded when the LOST packet was generated.
    \begin{itemize}
        \item Bit[4:3]:
        \begin{itemize}
            \item 0: HP CPU packets were not discarded
            \item 3: HP CPU packets were discarded
            \item Other values: reserved
        \end{itemize}
        \item Bit[5]:
        \begin{itemize}
            \item 0: DMA packets were not discarded
            \item 1: DMA packets were discarded
        \end{itemize}
        \item Bit[6]: Reserved
    \end{itemize}
    \item \textbf{anchored} - indicates the location of the 32 bits in the data packet:
    \begin{itemize}
        \item 0: Lower 32 bits
        \item 1: Higher 32 bits
    \end{itemize}
\end{itemize}

\begin{longtable}{|l| m{8cm} |}
    \caption{DMA Access Source}\label{tab:assistdbg-dma-source}
    \\\hline
    \rowcolor{lightgray}
    \textbf{Value} & \textbf{Source} \\\hline
    \endfirsthead \hline \rowcolor{lightgray}
    \textbf{Value} & \textbf{Source} \\\hline
    \endhead \endfoot \hline \endlastfoot
    0 & HP CPU \\\hline
    1 & Reserved \\\hline
    2 & Reserved \\\hline
    3 & Reserved \\\hline
    4 & Reserved \\\hline
    5 & MEM\_MONITOR \\\hline
    6 & TRACE \\\hline 
    7 & Reserved \\\hline
    8 & PSRAM\_MEM\_MONITOR \\\hline
    9 \verb+~+ 15 & Reserved \\\hline
    16 \verb+~+ 31 & Refer to the peripheral corresponding to the value 0 \verb+~+ 15 in Chapter \ref{mod:dma} \textit{\nameref{mod:dma}} > Table \ref{tab:gdma-ahb-peri-sel}. For example, a value of 16 matches the peripheral corresponding to 0 in this table, and a value of 17 matches the peripheral corresponding to 1. \\\hline
\end{longtable}

The module's internal buffer is 32 bits wide. If the HP CPU and DMA bus access logging are enabled simultaneously and data is recorded at the same time, the HP CPU packets are cached in the buffer first, followed by the DMA packets. The Debug Assistant module automatically fetches the buffered data and stores it in the specified memory space with a 32-bit data width.

In loop mode, data looping multiple times in the storage memory may cause residual data, which can interfere with packet parsing. For example, the lower 32 bits of an HP CPU packet may overwrite, leaving higher 32-bit residual data. Therefore, users must filter out possible residual data when determining the starting position of the first valid packet. Read \hyperref[regdesc:MEMMONITORLOGMEMCURRENTADDRREG]{MEM\_MONITOR\_LOG\_MEM\_CURRENT\_ADDR\_REG} to identify the starting position of the packet, then check the anchored bit value of the packet. If it is 0, retain the data. If it is 1, discard it. 

The packet parsing process is described below.

\begin{itemize}
    \item Read \hyperref[fielddesc:MEMMONITORLOGMEMFULLFLAG]{MEM\_MONITOR\_LOG\_MEM\_FULL\_FLAG} to determine whether there is a data overflow.
    \begin{itemize}
        \item If no, the address space to read is \hyperref[regdesc:MEMMONITORLOGMEMSTARTREG]{MEM\_MONITOR\_LOG\_MEM\_START\_REG} \verb+~+ \hyperref[regdesc:MEMMONITORLOGMEMCURRENTADDRREG]{MEM\_MONITOR\_LOG\_MEM\_CURRENT\_ADDR\_REG} $-$ 4.
        \item If yes and loop mode is enabled, the address space is \hyperref[regdesc:MEMMONITORLOGMEMCURRENTADDRREG]{MEM\_MONITOR\_LOG\_MEM\_CURRENT\_ADDR\_REG} \verb+~+ \hyperref[regdesc:MEMMONITORLOGMEMENDREG]{MEM\_MONITOR\_LOG\_MEM\_END\_REG} and \hyperref[regdesc:MEMMONITORLOGMEMSTARTREG]{MEM\_MONITOR\_LOG\_MEM\_START\_REG} \verb+~+ \hyperref[regdesc:MEMMONITORLOGMEMCURRENTADDRREG]{MEM\_MONITOR\_LOG\_MEM\_CURRENT\_ADDR\_REG} $-$ 4.
        \item If yes and loop mode is not enabled, the address space is \hyperref[regdesc:MEMMONITORLOGMEMSTARTREG]{MEM\_MONITOR\_LOG\_MEM\_START\_REG} \verb+~+ \hyperref[regdesc:MEMMONITORLOGMEMENDREG]{MEM\_MONITOR\_LOG\_MEM\_END\_REG}.
    \end{itemize}
    \item Read and parse data from the starting address. Read 32 bits at a time.
\end{itemize}

After completing packet parsing, clear the \hyperref[fielddesc:MEMMONITORLOGMEMFULLFLAG]{MEM\_MONITOR\_LOG\_MEM\_FULL\_FLAG} flag bit by setting \hyperref[fielddesc:MEMMONITORCLRLOGMEMFULLFLAG]{MEM\_MONITOR\_CLR\_LOG\_MEM\_FULL\_FLAG} to 1.

\subsubsection{Bus Access Logging 2 Configuration}

Bus access logging 2 includes DMA bus access logging (external RAM monitoring). The configuration process is described below.

\begin{enumerate}
    \item Configure monitored address space: Set \hyperref[regdesc:MEMMONITORLOGMINREG]{MEM\_MONITOR\_LOG\_MIN\_REG} and \hyperref[regdesc:MEMMONITORLOGMAXREG]{MEM\_MONITOR\_LOG\_MAX\_REG} to specify the monitored address space, which should range from 0x4200\_0000 to 0x43FF\_FFFF.
    \item Configure the monitoring mode with \hyperref[fielddesc:MEMMONITORLOGMODE]{MEM\_MONITOR\_LOG\_MODE}:
        \begin{itemize}
            \item Write monitoring (detects bus write operations)
            \item Word monitoring (detects writes of a specific word)
            \item Halfword monitoring (detects writes of a specific halfword)
            \item Byte monitoring (detects writes of a specific byte)
        \end{itemize}
    \item Configure the specific values to monitor:
        \begin{itemize}
            \item In word monitoring mode, \hyperref[regdesc:MEMMONITORLOGCHECKDATAREG]{MEM\_MONITOR\_LOG\_CHECK\_DATA\_REG} specifies the monitored word.
            \item In halfword monitoring mode, \hyperref[regdesc:MEMMONITORLOGCHECKDATAREG]{MEM\_MONITOR\_LOG\_CHECK\_DATA\_REG}[15:0] specifies the monitored halfword.
            \item In byte monitoring mode, \hyperref[regdesc:MEMMONITORLOGCHECKDATAREG]{MEM\_MONITOR\_LOG\_CHECK\_DATA\_REG}[7:0] specifies the monitored byte.
            \item Use \hyperref[regdesc:MEMMONITORLOGDATAMASKREG]{MEM\_MONITOR\_LOG\_DATA\_MASK\_REG} to mask the byte specified in \hyperref[regdesc:MEMMONITORLOGCHECKDATAREG]{MEM\_MONITOR\_LOG\_CHECK\_DATA\_REG}. A masked byte can be any value. For example, if \hyperref[regdesc:MEMMONITORLOGCHECKDATAREG]{MEM\_MONITOR\_LOG\_CHECK\_DATA\_REG} is set to 0x01020304 and \hyperref[regdesc:MEMMONITORLOGDATAMASKREG]{MEM\_MONITOR\_LOG\_DATA\_MASK\_REG} is set to 0x1, then any writes matching the 0x010203XX pattern will be recorded.
        \end{itemize}
    \item Configure the storage space for recorded data:
        \begin{itemize}
            \item \hyperref[regdesc:MEMMONITORLOGMEMSTARTREG]{MEM\_MONITOR\_LOG\_MEM\_START\_REG} and \hyperref[regdesc:MEMMONITORLOGMEMENDREG]{MEM\_MONITOR\_LOG\_MEM\_END\_REG} specify the storage space for recorded data, which must be in the range of 0x4080\_0000 to 0x4084\_FFFF.
            \item Set \hyperref[regdesc:MEMMONITORLOGMEMADDRUPDATEREG]{MEM\_MONITOR\_LOG\_MEM\_ADDR\_UPDATE\_REG} to update the value in \hyperref[regdesc:MEMMONITORLOGMEMCURRENTADDRREG]{MEM\_MONITOR\_LOG\_MEM\_CURRENT\_ADDR\_REG} to \hyperref[regdesc:MEMMONITORLOGMEMSTARTREG]{MEM\_MONITOR\_LOG\_MEM\_START\_REG}.
            \item Configure the Debug Assistant module's permission to access the HP SRAM. The Debug Assistant module can access the HP SRAM only when access permission is enabled. For more information, refer to Chapter \ref{mod:permctrl} \textit{\nameref{mod:permctrl}}.
        \end{itemize}
    \item Configure the writing mode for recorded data in loop mode or non-loop mode:
        \begin{itemize}
            \item In loop mode, writing to the specified address space occurs in loops. When writing reaches the end address, it returns to the starting address and continues, overwriting previously recorded data. Set \hyperref[fielddesc:MEMMONITORLOGMEMLOOPENABLE]{MEM\_MONITOR\_LOG\_MEM\_LOOP\_ENABLE} to enable loop mode.
            \item In non-loop mode, when writing reaches the end address, it stops and discards the remaining data, without overwriting previously recorded data. Clear \hyperref[fielddesc:MEMMONITORLOGMEMLOOPENABLE]{MEM\_MONITOR\_LOG\_MEM\_LOOP\_ENABLE} to enable non-loop mode.
            \item See examples in Section \ref{sec:hp-cpu-access-mon} > step \ref{other:dbgassist-item-log-config-step5}.
        \end{itemize}
    \item Set \hyperref[fielddesc:MEMMONITORLOGDMA0ENA]{MEM\_MONITOR\_LOG\_DMA\_0\_ENA} to enable DMA\_0 channel groups bus access logging.
\end{enumerate}

The Debug Assistant module first writes the DMA recorded data to the internal buffer B, then fetches the data from the buffer and writes it to the configured memory space. When monitored behaviors are triggered continuously, the generated recording packets may fill the buffer, preventing it from accepting new packets. In such cases, the module discards incoming packets and buffers a DMA LOST packet to prevent the buffer from reaching full capacity. The DMA LOST packet contains information about the bus type of the discarded packets, but not the number of discarded packets.

When bus access logging is finished, read the recorded data from memory for decoding. The recorded data is in two packet formats: DMA\_0 packet (corresponding to DMA\_0 Data bus) and DMA LOST packet. The packet formats are shown in Table \ref{tab:dbgassist-dma0-package} and \ref{tab:dbgassist-dma-lost-package}.

\begin{longtable}[c]{|L{2cm}|L{2cm}|L{2cm}|L{2cm}|}
    \caption{DMA\_0 Packet Format}\label{tab:dbgassist-dma0-package}
    \\\hline
    \rowcolor{lightgray}
    \textbf{Bit[31:7]} & \textbf{Bit[6:2]} & \textbf{Bit[1]} & \textbf{Bit[0]} \\\hline
    addr\_offset & dma\_source & format & anchored(0)  \\\hline
\end{longtable}
    
\begin{longtable}[c]{|L{2cm}|L{2cm}|L{2cm}|L{2cm}|}
    \caption{DMA LOST Packet Format}\label{tab:dbgassist-dma-lost-package}
    \\\hline
    \rowcolor{lightgray}
    \textbf{Bit[31:3]} & \textbf{Bit[2]} & \textbf{Bit[1]} & \textbf{Bit[0]} \\\hline
    reserved & lost\_source & format & anchored(0) \\\hline
\end{longtable}

The data packet formats show that the DMA\_0 and DMA LOST packets are both 32 bits. These packets contain the following fields:

\begin{itemize}
    \item \textbf{format} – indicates the source of the packet type:
    \begin{itemize}
        \item 0: DMA\_0 packet
        \item 1: DMA LOST packet
    \end{itemize}
    \item \textbf{addr\_offset} – indicates the address offset of a write operation. Actual address = addr\_offset + \{\hyperref[regdesc:MEMMONITORLOGMINREG]{MEM\_MONITOR\_LOG\_MIN\_REG}[31:2], 2'b0\}.
    \item \textbf{dma\_source} – the source of DMA access. Details are in Table \ref{tab:assistdbg-dma-source}.
    \item \textbf{lost\_source} – indicates which packets were discarded when the DMA LOST packet was generated:
    \begin{itemize}
        \item 0: DMA\_0 packets were not discarded
        \item 1: DMA\_0 packets were discarded
    \end{itemize}
    \item \textbf{anchored} – indicates the location of the 32 bits in the data packet:
    \begin{itemize}
        \item 0: Lower 32 bits
        \item 1: Higher 32 bits
    \end{itemize}
\end{itemize}

The internal buffer of the module is 32 bits wide. When the bus access logging of the DMA\_0 channel groups is enabled and the record data is generated, the DMA\_0 data packets are buffered. The Debug Assistant module automatically fetches the buffered data and stores it in 32-bit width in the specified memory space.

The packet parsing process is described below.

\begin{itemize}
    \item Read \hyperref[fielddesc:MEMMONITORLOGMEMFULLFLAG]{MEM\_MONITOR\_LOG\_MEM\_FULL\_FLAG} to determine whether there is a data overflow.
    \begin{itemize} 
        \item If no, the address space to read is \hyperref[regdesc:MEMMONITORLOGMEMSTARTREG]{MEM\_MONITOR\_LOG\_MEM\_START\_REG}  \verb+~+ \hyperref[regdesc:MEMMONITORLOGMEMCURRENTADDRREG]{MEM\_MONITOR\_LOG\_MEM\_CURRENT\_ADDR\_REG} $-$ 4.
        \item If yes and loop mode is enabled, the address space is \hyperref[regdesc:MEMMONITORLOGMEMCURRENTADDRREG]{MEM\_MONITOR\_LOG\_MEM\_CURRENT\_ADDR\_REG} \verb+~+ \hyperref[regdesc:MEMMONITORLOGMEMENDREG]{MEM\_MONITOR\_LOG\_MEM\_END\_REG} and \hyperref[regdesc:MEMMONITORLOGMEMSTARTREG]{MEM\_MONITOR\_LOG\_MEM\_START\_REG} \verb+~+ \hyperref[regdesc:MEMMONITORLOGMEMCURRENTADDRREG]{MEM\_MONITOR\_LOG\_MEM\_CURRENT\_ADDR\_REG} $-$ 4.
        \item If yes and loop mode is not enabled, the address space is \hyperref[regdesc:MEMMONITORLOGMEMSTARTREG]{MEM\_MONITOR\_LOG\_MEM\_START\_REG} \verb+~+ \hyperref[regdesc:MEMMONITORLOGMEMENDREG]{MEM\_MONITOR\_LOG\_MEM\_END\_REG}.
    \end{itemize}
    \item Read and parse data from the starting address. Read 32 bits at a time.
\end{itemize}

After packet parsing is complete, clear the \hyperref[fielddesc:MEMMONITORLOGMEMFULLFLAG]{MEM\_MONITOR\_LOG\_MEM\_FULL\_FLAG} flag bit by setting \hyperref[fielddesc:MEMMONITORCLRLOGMEMFULLFLAG]{MEM\_MONITOR\_CLR\_LOG\_MEM\_FULL\_FLAG} to 1.


%%%%%%%%%%%%%%%%%%%%%%%%
%%% Register Summary %%%
%%%%%%%%%%%%%%%%%%%%%%%%

\newpage
\begin{landscape}
\hypertarget{dbgassist-reg-summ}{}
\section{Register Summary}
\label{sec:dbgassist-reg-summ}

The addresses of \textbf{bus logging configuration registers} in Section \ref{sec:dbgassist-mem-monitor-reg-summ} are relative to the \textbf{TCM\_MEM\_MONITOR} base address and the \textbf{PSRAM\_MEM\_MONITOR}. The addresses of other registers in Section \ref{sec:dbgassist-bus-monitor-reg-summ} are relative to the \textbf{BUS\_MONITOR} base address. All base addresses are provided in Table \ref{tab:sysmem-base-address} in Chapter \ref{mod:sysmem} \textit{\nameref{mod:sysmem}}.

The abbreviations given in Column \textbf{Access} are explained in Section \hyperref[glossary-access-types]{\textit{Access Types for Registers}}.

\subsection{Bus Logging Configuration Registers}
\label{sec:dbgassist-mem-monitor-reg-summ}

\subfile{\modulefiles/19-DBGASSIST-mem-monitor-reg-summ__EN}

\subsection{SOther Registers}
\label{sec:dbgassist-bus-monitor-reg-summ}

\subfile{\modulefiles/19-DBGASSIST-bus-monitor-reg-summ__EN}
\end{landscape}

%%%%%%%%%%%%%%%%%
%%% Registers %%%
%%%%%%%%%%%%%%%%%

% >>> Update the sample text below and insert
%      the info on registers into the subfile <<<

\newpage
\section{Registers}

The addresses of \textbf{bus logging configuration registers} in Section \ref{sec:dbgassist-mem-monitor-reg} are relative to the \textbf{TCM\_MEM\_MONITOR} base address and the \textbf{PSRAM\_MEM\_MONITOR}. The addresses of other registers in Section \ref{sec:dbgassist-bus-monitor-reg} are relative to the \textbf{BUS\_MONITOR} base address. All base addresses are provided in Table \ref{tab:sysmem-base-address} in Chapter \ref{mod:sysmem} \textit{\nameref{mod:sysmem}}.

For how to program reserved fields, please refer to Section \hyperref[programming-reserved-field]{\textit{Programming Reserved Register Field}}.

\subsection{Bus Logging Configuration Registers}
\label{sec:dbgassist-mem-monitor-reg}

\subfile{\modulefiles/19-DBGASSIST-mem-monitor-reg__EN}

\subsection{Other Registers}
\label{sec:dbgassist-bus-monitor-reg}

\subfile{\modulefiles/19-DBGASSIST-bus-monitor-reg__EN}

\end{document}
